\documentclass{article}
\usepackage[UTF8]{ctex}
\usepackage{amsmath,amssymb,amsfonts}
\usepackage{enumitem}
\usepackage{geometry}
\geometry{a4paper, margin=1in}
% === 错题本排版优化（不参与题目标识） ===
\setlength{\parindent}{0pt}
\setlength{\parskip}{0.35em plus 0.10em minus 0.08em}
\setlength{\emergencystretch}{3em}
\setlist{itemsep=0.25em, topsep=0.40em, parsep=0pt}
\allowdisplaybreaks[2]
% === 错题本排版优化结束 ===
\begin{document}
\section*{华中科技大学2026年高等代数真题}
\section*{125．华中科技大学2026年高等代数试卷}
\begin{itemize}

\subsection{题目1}
  \item[1．] 计算行列式
\[
\left|\begin{array}{ccccc}
1 & x_{1} & \cdots & x_{1}^{n-2} & \left(x_{2}+x_{3}+\cdots+x_{n}\right)^{n-1} \\
1 & x_{2} & \cdots & x_{2}^{n-2} & \left(x_{1}+x_{3}+\cdots+x_{n}\right)^{n-1} \\
\vdots & \vdots & & \vdots & \vdots \\
1 & x_{n} & \cdots & x_{n}^{n-2} & \left(x_{1}+x_{2}+\cdots+x_{n-1}\right)^{n-1}
\end{array}\right|
\]

\subsection{题目2}
  \item[2．] 设 \(V_{1}=L\left(\alpha_{1}, \alpha_{2}, \alpha_{3}\right), V_{2}=L\left(\beta_{1}, \beta_{2}\right)\) ，求 \(V_{1}+V_{2}, V_{1} \cap V_{2}\) 的基与维数（具体数据忘了）。
\end{itemize}

\subsection{题目3}
3．\(C \in M_{m \times n}(\mathbb{C}), r(C)=r\) ，且 \(A C=C B\) ，证明：\(A, B\) 至少有 \(r\) 个公共的特征值（计重数）．

\subsection{题目4}
4．有限维线性空间 \(V\) 有2026个子空间 \(W_{1}, W_{2}, \cdots, W_{2026}\) ，其中

\[
\operatorname{dim} W_{i}=2026(i=1,2, \cdots, 2026), \operatorname{dim}\left(W_{i} \cap W_{j}\right)=2025(i \neq j) .
\]

证明下列条件之一成立：

\begin{itemize}
  \item[（a）] 存在 \(W\) 为 \(V\) 的子空间，且 \(\operatorname{dim} W=2025, W \subset W_{i}(i=1,2, \cdots, 2026)\) 。
  \item[（b）] 存在 \(U\) 为 \(V\) 的子空间，且 \(\operatorname{dim} U=2027, W_{i} \subset U(i=1,2, \cdots, 2026)\) ．
\end{itemize}

\subsection{题目5}
5．设 \(A \in M_{n \times n}(\mathbb{C})\) ．证明：

\begin{itemize}
  \item[（1）] 存在 \(\delta>0\) ，使得 \(\lambda \in \mathbb{C} \backslash\{0\}\) 且 \(|\lambda|<\delta\) 时，\(\lambda E+A\) 可逆．
  \item[（2）] \(\lim _{\lambda \rightarrow 0} A(\lambda E+A)^{-1}\) 存在的充分必要条件是 \(r(A)=r\left(A^{2}\right)\) ．
\end{itemize}

\subsection{题目6}
6．（可能有误）设 \(A\) 为三阶可逆实对称矩阵，证明下面两个命题等价：

\begin{itemize}
  \item[（a）] \(\operatorname{tr}(A) \neq 0\) ．
  \item[（b）] 对任意三阶实反对称矩阵 \(S\) ，存在三阶实反对称矩阵 \(B\) ，使得 \(A B+B A=S\) ．
\end{itemize}

\subsection{题目7}
7．\(A_{n \times n}\) 为正定矩阵，\(B_{n \times m}\) 满足 \(r(B)=r\) ，求矩阵 \(M=\left(\begin{array}{cc}A & B \\ B^{\mathrm{T}} & O\end{array}\right)\) 的正负惯性指数。

\subsection{题目8}
8．（可能有误）\(A, B\) 为2026阶反对称矩阵，证明：

\begin{itemize}
  \item[（1）] \(r(A)\) 为偶数．
  \item[（2）] \(A B\) 的特征多项式至少有一个二重根．
\end{itemize}
\end{document}
